\section*{Appendix: machine-checkable $E_8=(D_5\oplus A_3)+\mu_4$ construction}
\addcontentsline{toc}{section}{Appendix: explicit $E_8$ construction}
% =====================================================================
The lift is reproducible by a short script (all entries below verified exactly).

\smallskip
\noindent\textbf{Cartan matrices and discriminant groups.}
\[
  C(D_5)=\begin{pmatrix}2&-1&0&0&0\\-1&2&-1&0&0\\0&-1&2&-1&-1\\0&0&-1&2&0\\0&0&-1&0&2\end{pmatrix},\quad
  C(A_3)=\begin{pmatrix}2&-1&0\\-1&2&-1\\0&-1&2\end{pmatrix};\quad
  \det C(D_5)=\det C(A_3)=4 .
\]
Smith normal forms give $\operatorname{disc}(D_5)=\operatorname{disc}(A_3)=\mathbb Z_4$; the glue index is
$[\,E_8:D_5\oplus A_3\,]=\sqrt{4\cdot4/1}=4=|\mu_4|$, with discriminant-form norms $q(D_5)=\tfrac54$,
$q(A_3)=\tfrac34$, $q(D_5)+q(A_3)=2\equiv0\pmod 2$ (even-glue). Writing $\mu:=|\mu_4|=4$ the norms are
$q(D_5)=\tfrac{\mu+1}{\mu}$, $q(A_3)=\tfrac{\mu-1}{\mu}$, so $q(D_5)q(A_3)=1-\tfrac1{\mu^2}=\tfrac{15}{16}$ and
\begin{equation}
  \boxed{\ \dim S^+\bigl(1-q(D_5)q(A_3)\bigr)=16\cdot\tfrac1{16}=1\ }\qquad\mE ,
\end{equation}
i.e.\ the glue product fills all $15=\dim S^+{-}1$ non-trivial half-spin states and the defect is exactly
the singleton (the $\nu^c$/Higgs closure). And $E_8$ has a tidy $\mu_4$-cardinal form,
$\dim E_8=\binom{|\mu_4|^2}{2}+2|\mu_4|^3=\binom{16}{2}+2\cdot64=120+128=248$ \mC.

\smallskip
\noindent\textbf{Explicit roots.} Split $\mathbb R^8=\mathbb R^5_{(D_5)}\oplus\mathbb R^3_{(A_3)}$. The
$240$ roots are
\[
  \{\pm e_i\pm e_j:1\le i<j\le 8\}\ (112)\ \cup\ \{\tfrac12(\pm1)^8:\ \text{even \# of }-\}\ (128),
\]
all of norm $2$. Splitting each by $(\|\cdot\|^2_{D_5},\|\cdot\|^2_{A_3})$ gives the branching blocks
\[
  (2,0)\!:40\ [D_5],\quad (0,2)\!:12\ [A_3],\quad (1,1)\!:60\ [(10,6)],\quad
  (\tfrac54,\tfrac34)\!:128\ [(16,4)\oplus(\overline{16},\overline4)] ,
\]
so $\mathfrak e_8=(45,1)\oplus(1,15)\oplus(10,6)\oplus(16,4)\oplus(\overline{16},\overline4)$,
$248=45{+}15{+}60{+}64{+}64$, $240=40{+}12{+}60{+}64{+}64$.

\smallskip
\noindent\textbf{Simple-root Gram (unimodular check).} With the Bourbaki simple roots
$\alpha_1=\tfrac12(e_1{-}e_2{-}\dots{-}e_7{+}e_8)$, $\alpha_2=e_1{+}e_2$, and
$\alpha_{k}=e_{k-1}{-}e_{k-2}$ for $k=3,\dots,8$ (so $\alpha_3=e_2{-}e_1,\ \alpha_4=e_3{-}e_2,\ \dots,\ \alpha_8=e_7{-}e_6$),
the Gram matrix $G_{ij}=\langle\alpha_i,\alpha_j\rangle$ is the $E_8$ Cartan matrix with $\det G=1$ and
all $G_{ii}=2$; every one of the $240$ roots is an integer combination of the $\alpha_i$. Hence the
assembled lattice is $E_8$, not merely a $240$-count. \mE\ (The shifted index $e_{k-2}{-}e_{k-3}$ would
duplicate $\alpha_3$ and make $G$ singular --- use $e_{k-1}{-}e_{k-2}$.)

\clearpage
\part{Appendix: ablation of the EM fixed point}
\label{sec:ablation}

The single sharpest reviewer test is: \emph{was the closure $F_{U(1)}(\alpha)=0$ engineered so that
$\alpha$ comes out?} The honest answer is an ablation --- vary each input and see which actually move the
root. Three one-parameter scans (all others held at their compiler values) separate the
\emph{load-bearing} integers from the \emph{fine} last-digit selector.

\begin{figure}[ht]
\centering
\includegraphics[width=0.92\linewidth]{figures/alpha_ablation.pdf}
\caption{EM-closure ablation (data plot, \texttt{verification/make\_figures.py}). \textbf{(a)} $\ainv$ vs
the integer budget $M$: only $M=41$ lands on CODATA-2022. \textbf{(b)} $\ainv$ vs the seam denominator
$N$ in $\cthree=1/(N\pi)$: only $N=8$ lands. Both integers are fixed elsewhere (flavor budget; rank/$h(D_5)$),
so the equation is not freely tunable to $137.036$.}
\label{fig:ablation}
\end{figure}

% --- ablation 1 ---
\noindent\textbf{(A1) The transport budget $M=\sum L+N_\Phi$} (exponent held at $-\tfrac54$):
\begin{center}\footnotesize\renewcommand{\arraystretch}{1.15}
\begin{tabular}{@{}cccccccc@{}}
\toprule
$M$ & $38$ & $39$ & $40$ & $\mathbf{41}$ & $42$ & $43$ & $44$ \\
\midrule
$\ainv$ & $140.530$ & $139.326$ & $138.162$ & $\mathbf{137.0359992}$ & $135.946$ & $134.890$ & $133.866$ \\
\bottomrule
\end{tabular}
\end{center}
Each unit of $M$ shifts $\ainv$ by $\approx1.09$, so \emph{only} $M=41$ lands on $137.036$. This is a hard
lock --- and $41=10b_1=\sum L+N_\Phi=40+1$ is fixed by the compiler (flavor budget $40$ $+$ one Higgs).
\textbf{Load-bearing.}

% --- ablation 2 ---
\noindent\textbf{(A2) The seam ``8'' in $\cthree=1/(N\pi)$} (here $\dtop=48\cthree^4$ tracks $\cthree$):
\begin{center}\footnotesize\renewcommand{\arraystretch}{1.15}
\begin{tabular}{@{}cccccc@{}}
\toprule
$N$ & $6$ & $7$ & $\mathbf{8}$ & $9$ & $10$ \\
\midrule
$\ainv$ & $76.97$ & $104.85$ & $\mathbf{137.036}$ & $173.53$ & $214.34$ \\
\bottomrule
\end{tabular}
\end{center}
Only $N=8$ --- i.e.\ $\cthree=1/(8\pi)=1/(\operatorname{rank}E_8\cdot\pi)$ --- reproduces $137.036$. The
``8'' is the same one fixed by the bootstrap ($h(D_5)=\operatorname{rank}E_8=\varphi(30)$) and the
Gauss--Bonnet hardening. \textbf{Load-bearing.}

% --- ablation 3 ---
\noindent\textbf{(A3) The carrier discriminant exponent} (at $M=41$, $\cthree=1/(8\pi)$):
\begin{center}\footnotesize\renewcommand{\arraystretch}{1.15}
\begin{tabular}{@{}cccccc@{}}
\toprule
exponent & $-\tfrac14$ & $-\tfrac34$ & $\mathbf{-\tfrac54}$ & $-\tfrac74$ & $-\tfrac94$ \\
\midrule
$\ainv$ & $137.035995$ & $137.035997$ & $\mathbf{137.0359992}$ & $137.036001$ & $137.036003$ \\
\bottomrule
\end{tabular}
\end{center}
The exponent moves only the \emph{seventh} digit; all candidates agree to six. Here $-\tfrac54=q(D_5)$ is
selected, but this is a \emph{fine} selector, not a hard lock. \textbf{Fine (last-digit).}

\begin{keybox}[Ablation verdict --- honest]
The EM fixed point is a \textbf{hard test of two integers} --- the budget $M=41$ and the seam denominator
$8=\operatorname{rank}E_8$, each decisive (a single-unit change misses by $\sim1$ in $\ainv$) --- and a
\textbf{fine test of one norm}, the carrier discriminant exponent $-\tfrac54=q(D_5)$ (seventh digit). The
equation is therefore \emph{not} freely tunable to $137.036$: the two integers that hit it are exactly the
compiler/bootstrap numbers, fixed elsewhere. What it does \emph{not} do is pin the carrier norm strongly
--- that is honest, and it is the right level of claim. \mE
\end{keybox}

% =====================================================================
\section*{Appendix: computational verification}
\addcontentsline{toc}{section}{Appendix: computational verification}
% =====================================================================
The load-bearing identities of this document (marked \mE, \mE, \mE) are re-derived from
$\{\cthree,\gcar\}$ alone and machine-checked by the suite in the repository folder
\texttt{verification/}; the inline tags \texttt{[verification/$\dots$]} point to the exact script.
Relevant here:
\begin{itemize}[leftmargin=1.4em,itemsep=1pt]
\item \texttt{v1\_e8\_glue.py} --- the $\mu_4$ glue ($q(D_5){+}q(A_3){=}2$, $\operatorname{disc}{=}\Z_4$, $240,248$).
\item \texttt{v2\_carrier\_pascal.py} --- $\gcar{=}5$, $16{=}1{+}5{+}10$, $\Nfam,\Oadm,b_1$.
\item \texttt{v3\_em\_alpha.py} --- the EM fixed point $\ainv{=}137.0359992168$ (existence/uniqueness, ablation).
\item \texttt{v4\_flavor\_matrix.py} --- $R$ with $\det{=}8$, minors $\{2,3,5\}$, $\sum L{=}40$.
\item \texttt{v7\_gravity\_cosmo.py} --- scalaron $\cthree^7$, $\Lambda$/$A_s$ readouts.
\end{itemize}
Run \texttt{python verification/run\_all.py} (needs \texttt{mpmath, numpy, sympy}); it exits $0$ iff all
checks pass. Full map: \texttt{verification/README.md}.

\input{tex-artefacts/tfpt_references}

\end{document}
